\chapter{پیش‌درآمد: زبان و ابزارهای پایه}
\label{ch:prelim}

جبر خطی درباره‌ی \textbf{بردارها}، \textbf{ماتریس‌ها} و \textbf{دستگاه‌های معادلات خطی} است؛ اما برای این‌که بتوانیم دقیق حرف بزنیم و اثبات بنویسیم، به چند ابزار ساده نیاز داریم. این فصل کوتاه همین ابزارها را مرور می‌کند. اگر با مجموعه‌ها، نماد سیگما و اعداد مختلط راحت هستید، می‌توانید مستقیم به فصل ۱ بروید و هر وقت لازم شد به اینجا برگردید.

\section{مجموعه‌ها و نمادهای منطقی}

\subsection{مجموعه}
\term{مجموعه}{Set} یعنی گردایه‌ای از اشیا که به آن‌ها \term{عضو}{Element} می‌گوییم. چند نماد پرکاربرد:

\begin{center}
\begin{tabular}{>{\centering\arraybackslash}m{3.2cm} m{9cm}}
\toprule
\textbf{نماد} & \textbf{معنی} \\
\midrule
$x\in A$ & $x$ عضو مجموعه‌ی $A$ است. \\
$x\notin A$ & $x$ عضو $A$ نیست. \\
$A\subseteq B$ & $A$ زیرمجموعه‌ی $B$ است: هر عضو $A$ عضو $B$ هم هست. \\
$A\cap B$ & اشتراک: اعضایی که هم در $A$ و هم در $B$ هستند. \\
$A\cup B$ & اجتماع: اعضایی که دست‌کم در یکی از $A$ یا $B$ هستند. \\
$A\setminus B$ & تفاضل: اعضای $A$ که در $B$ نیستند. \\
$\emptyset$ & مجموعه‌ی تهی (بدون عضو). \\
$A\times B$ & ضرب دکارتی: همه‌ی زوج‌های مرتب $(a,b)$ با $a\in A$ و $b\in B$. \\
$\{x\in A : P(x)\}$ & اعضایی از $A$ که ویژگی $P$ را دارند. \\
\bottomrule
\end{tabular}
\end{center}

مثلاً $\{x\in\R : x^2=1\}=\{-1,1\}$.

\begin{note}[اثبات تساوی دو مجموعه]
برای این‌که ثابت کنیم $A=B$، معمولاً دو کار جدا انجام می‌دهیم: نشان می‌دهیم $A\subseteq B$ و سپس $B\subseteq A$. این روش را در فصل‌های زیرفضا و رتبه بارها به کار می‌بریم.
\end{note}

\subsection{مجموعه‌های عددی}
\begin{itemize}
  \item $\N=\{1,2,3,\dots\}$: اعداد طبیعی.
  \item $\Z=\{\dots,-2,-1,0,1,2,\dots\}$: اعداد صحیح.
  \item $\Q$: اعداد گویا (کسرهایی مثل $\frac{p}{q}$ با $p,q\in\Z$ و $q\neq 0$).
  \item $\R$: اعداد حقیقی.
  \item $\C$: اعداد مختلط (بخش \ref{sec:complex}).
\end{itemize}
همچنین $\R^n$ مجموعه‌ی همه‌ی $n$تایی‌های مرتب از اعداد حقیقی است؛ مثلاً $(1,-2,5)\in\R^3$.

\subsection{نمادهای منطقی}
\begin{center}
\begin{tabular}{>{\centering\arraybackslash}m{2.5cm} m{10cm}}
\toprule
\textbf{نماد} & \textbf{معنی} \\
\midrule
$\forall$ & «برای هر». مثلاً $\forall x\in\R:\ x^2\ge 0$. \\
$\exists$ & «وجود دارد». مثلاً $\exists x\in\R:\ x^2=2$. \\
$P\Rightarrow Q$ & «اگر $P$ آن‌گاه $Q$» (از $P$ نتیجه می‌شود $Q$). \\
$P\Leftrightarrow Q$ & «$P$ اگر و تنها اگر $Q$»؛ یعنی هم $P\Rightarrow Q$ و هم $Q\Rightarrow P$. \\
$:=$ & «تعریف می‌کنیم». مثلاً $f(x):=x^2$. \\
\bottomrule
\end{tabular}
\end{center}

\section{چند روش اثبات}
در این جزوه اثبات‌های زیادی می‌بینیم. خوشبختانه بیشترشان از چند الگوی ساده پیروی می‌کنند:

\begin{itemize}
  \item \textbf{اثبات مستقیم:} از فرض شروع می‌کنیم و قدم‌به‌قدم به حکم می‌رسیم.
  \item \textbf{اثبات «اگر و تنها اگر»:} حکمی مثل $P\Leftrightarrow Q$ دو بخش دارد. یک بار فرض می‌کنیم $P$ درست است و $Q$ را نتیجه می‌گیریم (جهت رفت) و یک بار برعکس (جهت برگشت).
  \item \textbf{برهان خلف} \enterm{Proof by Contradiction}: فرض می‌کنیم حکم \emph{نادرست} است و نشان می‌دهیم این فرض به یک تناقض می‌رسد؛ پس حکم باید درست باشد.
  \item \textbf{مثال نقض} \enterm{Counterexample}: برای رد کردن یک ادعای کلی («برای هر …»)، کافی است \emph{یک} مثال پیدا کنیم که ادعا در آن غلط باشد.
  \item \textbf{استقرای ریاضی} \enterm{Induction}: برای حکمی درباره‌ی همه‌ی اعداد طبیعی $n$، آن را برای $n=1$ ثابت می‌کنیم و نشان می‌دهیم اگر برای $n$ درست باشد برای $n+1$ هم درست است.
\end{itemize}

\section{نماد سیگما}
برای نوشتن جمع‌های طولانی از نماد $\sum$ استفاده می‌کنیم:
\[
\sum_{i=1}^{n} a_i := a_1+a_2+\dots+a_n .
\]
اندیس $i$ فقط یک «شمارنده» است و اسمش اهمیتی ندارد: $\sum_{i=1}^n a_i=\sum_{k=1}^n a_k$.

سه قاعده‌ی مهم که بارها از آن‌ها استفاده می‌کنیم:
\begin{align*}
\sum_{i=1}^{n} c\,a_i &= c\sum_{i=1}^{n} a_i, &
\sum_{i=1}^{n} (a_i+b_i) &= \sum_{i=1}^{n} a_i+\sum_{i=1}^{n} b_i, &
\sum_{i=1}^{m}\sum_{j=1}^{n} a_{ij} &= \sum_{j=1}^{n}\sum_{i=1}^{m} a_{ij}.
\end{align*}
قاعده‌ی سوم می‌گوید در یک جمع دوگانه می‌توانیم ترتیب جمع‌زدن را عوض کنیم؛ مثل این‌که اعداد یک جدول را یک بار سطر به سطر و یک بار ستون به ستون جمع بزنیم. نتیجه یکی است.

\section{تابع}
یک \term{تابع}{Function} $f:A\to B$ به هر عضو $a\in A$ دقیقاً یک عضو $f(a)\in B$ نسبت می‌دهد. به $A$ \term{دامنه}{Domain} و به $B$ \term{هم‌دامنه}{Codomain} می‌گوییم.
\begin{itemize}
  \item دو تابع $f,g:A\to B$ برابرند اگر برای \emph{هر} $a\in A$ داشته باشیم $f(a)=g(a)$.
  \item ترکیب دو تابع: $(g\circ f)(a):=g(f(a))$.
  \item تابع $g$ را \term{وارون}{Inverse} $f$ می‌نامیم اگر $g(f(a))=a$ و $f(g(b))=b$ برای همه‌ی $a,b$.
\end{itemize}

\section{اعداد مختلط}\label{sec:complex}
عدد مختلط عددی به شکل $z=a+bi$ است که $a,b\in\R$ و $i^2=-1$. به $a$ \textbf{بخش حقیقی} و به $b$ \textbf{بخش موهومی} می‌گوییم. جمع و ضرب طبق قواعد معمول جبر انجام می‌شود و هرجا $i^2$ دیدیم به‌جایش $-1$ می‌گذاریم:
\[
(1+2i)(3-i)=3-i+6i-2i^2=3+5i+2=5+5i .
\]

\begin{definition}[مزدوج و قدرمطلق]
برای $z=a+bi$:
\begin{itemize}
  \item \term{مزدوج}{Conjugate} $z$ عبارت است از $\conj{z}=a-bi$.
  \item \term{قدرمطلق}{Modulus} $z$ عبارت است از $\abs{z}=\sqrt{a^2+b^2}$.
\end{itemize}
\end{definition}

ویژگی‌هایی که در فصل ضرب داخلی به آن‌ها نیاز داریم (همه با محاسبه‌ی مستقیم ثابت می‌شوند):
\[
z\conj{z}=a^2+b^2=\abs{z}^2\ \ge 0,\qquad
\conj{z+w}=\conj{z}+\conj{w},\qquad
\conj{zw}=\conj{z}\,\conj{w},\qquad
\conj{\conj{z}}=z,
\]
و $z$ حقیقی است اگر و تنها اگر $\conj{z}=z$. همچنین $z+\conj{z}=2a=2\operatorname{Re}(z)$ و $\operatorname{Re}(z)\le\abs{z}$.

\section{میدان}\label{sec:field}
در جبر خطی با «عددها» کار می‌کنیم: آن‌ها را جمع، تفریق، ضرب و تقسیم می‌کنیم. مجموعه‌ای از عددها که این چهار عمل در آن بدون دردسر انجام شود را \textbf{میدان} می‌نامیم.

\begin{definition}[میدان]
یک \term{میدان}{Field} مجموعه‌ای مثل $\F$ است همراه با دو عمل جمع ($+$) و ضرب ($\cdot$) به طوری که برای همه‌ی $a,b,c\in\F$:
\begin{enumerate}
  \item \textbf{بسته بودن:} $a+b\in\F$ و $ab\in\F$.
  \item \textbf{جابه‌جایی:} $a+b=b+a$ و $ab=ba$.
  \item \textbf{شرکت‌پذیری:} $(a+b)+c=a+(b+c)$ و $(ab)c=a(bc)$.
  \item \textbf{توزیع‌پذیری:} $a(b+c)=ab+ac$.
  \item \textbf{عضو خنثی:} عضوهای متمایز $0$ و $1$ در $\F$ وجود دارند که $a+0=a$ و $1\cdot a=a$.
  \item \textbf{قرینه:} برای هر $a$ عضوی مثل $-a$ هست که $a+(-a)=0$.
  \item \textbf{وارون ضربی:} برای هر $a\neq 0$ عضوی مثل $a^{-1}$ هست که $a\,a^{-1}=1$.
\end{enumerate}
\end{definition}

\begin{example}[میدان‌ها و غیرمیدان‌ها]
کدام یک میدان است؟ $\Q$، $\R$، $\C$، $\Z$، $\N$.
\end{example}
\begin{solution}
$\Q$، $\R$ و $\C$ میدان هستند. $\Z$ میدان نیست چون مثلاً $2$ وارون ضربی در $\Z$ ندارد ($\frac12\notin\Z$). $\N$ هم میدان نیست؛ حتی عضو $0$ و قرینه ندارد.
\end{solution}

\begin{note}
در این درس، اگر چیزی نگوییم، میدان همان $\R$ است. هرجا لازم باشد ($\C$) را صریح می‌گوییم. هر جا $\F$ دیدید، می‌توانید آن را $\R$ یا $\C$ تصور کنید.
\end{note}

\section{حروف یونانی}
در ریاضی برای اسکالرها زیاد از حروف یونانی استفاده می‌شود. این جدول کمک می‌کند آن‌ها را درست بخوانید:
\begin{center}
\begin{tabular}{cl@{\qquad}cl@{\qquad}cl}
\toprule
$\alpha$ & \en{alpha} (آلفا) & $\beta$ & \en{beta} (بتا) & $\gamma$ & \en{gamma} (گاما) \\
$\delta$ & \en{delta} (دلتا) & $\varepsilon$ & \en{epsilon} (اپسیلون) & $\theta$ & \en{theta} (تتا) \\
$\lambda$ & \en{lambda} (لاندا) & $\mu$ & \en{mu} (میو) & $\pi$ & \en{pi} (پی) \\
$\sigma,\Sigma$ & \en{sigma} (سیگما) & $\varphi$ & \en{phi} (فی) & $\omega$ & \en{omega} (اُمگا) \\
\bottomrule
\end{tabular}
\end{center}

\begin{definition}[دلتای کرونکر]
نماد $\delta_{ij}$ را این‌طور تعریف می‌کنیم:
\[
\delta_{ij}=\begin{cases}1 & i=j\\ 0 & i\neq j\end{cases}
\]
مثلاً درایه‌ی $(i,j)$ ماتریس همانی $I$ برابر $\delta_{ij}$ است.
\end{definition}
